\section{Background and Motivation}

SAKI targets co-located \rocksdb instances whose background work shares finite
host resources. RocksDB exposes useful local controls, but not the host-wide
placement decision: how that fixed service budget should be placed across
instances as tenants' write pressure changes at different times.

\subsection{LSM Background Work}

The log-structured merge-tree makes foreground writes inexpensive by buffering
updates and later merging them into on-disk components~\cite{oneil1996lsm}.
Compaction rewrites sorted files, removes obsolete entries, controls level
sizes, and bounds read cost. This design gives high write throughput but turns
recent foreground writes into future CPU and I/O demand. That deferred demand
becomes a host-level problem when several engines share one finite background
budget; surveys and recent microbenchmarks emphasize compaction, amplification,
and buffering as central LSM tradeoffs~\cite{luo2020lsmsurvey,
kaushik2024lsmbuffer}.

RocksDB exposes per-instance background jobs, rate limiters, statistics,
properties, and event listeners~\cite{rocksdb}. SAKI uses them as the
measurement and actuation surface for the host-level placement problem.

\subsection{From Tenant-Local State to Host-Level Debt}

We use \emph{compaction debt} as an operational term for deferred background
work and its foreground consequences. Pending compaction bytes, L0 files,
flush/compaction events, write-tail latency, and completion gaps are observable
but incomplete symptoms that help identify tenants currently creating deferred
maintenance.

The cross-tenant aspect is the key. Under a fixed aggregate budget, equal
allocation leaves capacity assigned to quiet tenants whenever only a subset is
write-heavy. Frozen skew helps only while the workload resembles the phase it was
tuned for. Section~\ref{sec:evaluation} tests this mismatch between where the
budget sits and where debt is being created.

Compaction service has delayed effects: too little service may surface later as
L0 pressure, flush backlog, write slowdown, or tail latency; too much service
for a quiet tenant consumes capacity that could reduce risk elsewhere. The
evaluation therefore reports high-pressure latency and throughput together with
total throughput and collateral impact on lower-pressure tenants.

\subsection{Why Existing Directions Do Not Settle the Question}

Prior LSM work changes single-engine policies or storage layout: Monkey,
Dostoevsky, WiscKey, PebblesDB, and SILK improve what an engine does given an
assigned service budget~\cite{dayan2017monkey,dayan2018dostoevsky,
lu2016wisckey,raju2017pebblesdb,balmau2019silk}. SAKI asks where a fixed host
budget should be placed across independent engines.

CaaS-LSM treats compaction as a service and DFlush offloads flush work to
DPUs~\cite{yu2024caaslsm,ding2025dflush}. SAKI shares the resource-management
view, but targets co-located RocksDB instances on one host, standard runtime
controls, and a conserved aggregate budget. The next section formalizes that
host-level allocation question.
